\section{Технический проект}

\subsection{Общая характеристика программного решения задачи}

Разработанная программно-информационная система обнаружения дефектов дорожного покрытия представляет собой настольное приложение, предназначенное для автоматизированного анализа изображений дорожного полотна с использованием методов компьютерного зрения и нейросетевой обработки данных.

Основной задачей системы является выявление дефектов дорожного покрытия на входных изображениях, их локализация и визуальное отображение результатов для последующего анализа пользователем.

Программное решение реализовано в виде локального приложения, не требующего постоянного сетевого соединения. Все этапы обработки данных выполняются на стороне пользователя, что обеспечивает автономность работы системы и снижает зависимость от внешних сервисов.

Обработка изображений осуществляется последовательно: входные данные загружаются пользователем через графический интерфейс, после чего передаются в модуль нейросетевой детекции. Результаты анализа сохраняются и отображаются в виде графической и текстовой информации.

Архитектура системы ориентирована на модульный принцип построения, при котором функциональность разделена на независимые компоненты, взаимодействующие через центральный управляющий модуль. Такой подход обеспечивает расширяемость системы, упрощает сопровождение программного кода и позволяет в дальнейшем добавлять новые функциональные возможности без изменения существующей логики.

Особое внимание при проектировании системы уделено обеспечению высокой скорости обработки изображений, рациональному использованию вычислительных ресурсов и корректному управлению графическим интерфейсом пользователя.

\subsection{Обоснование выбора технологий проектирования}

При разработке программно-информационной системы обнаружения дефектов дорожного покрытия выбор технологического стека осуществлялся исходя из требований к производительности, удобству разработки, возможности интеграции современных методов компьютерного зрения, а также обеспечения удобного графического пользовательского интерфейса.

В результате анализа существующих инструментальных средств был выбран технологический стек, включающий язык программирования Python, графический фреймворк PySide6, библиотеку OpenCV, нейросетевую модель семейства YOLOv8 и библиотеку NumPy.

\subsubsection{Язык программирования Python}

В качестве основного языка разработки программной системы был выбран Python.

Выбор данного языка обусловлен спецификой решаемой задачи. Разрабатываемая система сочетает в себе графический пользовательский интерфейс, обработку изображений и применение нейросетевой модели обнаружения объектов. Python предоставляет единое программное окружение для реализации всех перечисленных компонентов без необходимости использования нескольких языков программирования.

Существенным преимуществом Python является наличие развитой экосистемы библиотек для компьютерного зрения и машинного обучения. Это позволило использовать готовые инструменты для работы с изображениями, выполнения нейросетевого анализа и построения графического интерфейса приложения.

Дополнительным преимуществом языка является высокая скорость разработки и простота модификации программного кода. Это особенно важно при экспериментальной настройке параметров обработки изображений и тестировании различных вариантов работы модели обнаружения.

Следует отметить, что Python уступает компилируемым языкам программирования по скорости выполнения отдельных вычислительных операций. Однако в рамках данного проекта этот недостаток не оказывает существенного влияния на производительность системы, поскольку основные ресурсоёмкие вычисления выполняются средствами библиотеки Ultralytics и задействуют оптимизированные низкоуровневые реализации.

\subsubsection{Фреймворк PySide6}

Для разработки графического интерфейса выбран фреймворк PySide6, представляющий собой официальный набор Python-привязок для библиотеки Qt.

При выборе альтернатив рассматривались следующие решения:

\begin{enumerate}
	\item Tkinter --- отличается простотой использования, однако обладает ограниченными возможностями построения современных интерфейсов.
	\item Kivy --- ориентирован преимущественно на мультитач-интерфейсы и мобильные приложения.
\end{enumerate}

Выбор PySide6 обусловлен его ориентацией на создание полнофункциональных настольных приложений и тесной интеграцией с языком Python. Фреймворк предоставляет развитый набор визуальных компонентов, систему обработки событий и средства управления состоянием интерфейса.

К основным возможностям PySide6 относятся:

\begin{itemize}
	\item создание оконных приложений со сложной иерархией элементов;
	\item событийно-ориентированная модель взаимодействия;
	\item механизм сигналов и слотов;
	\item встроенные средства хранения пользовательских настроек;
	\item поддержка стилизации интерфейса;
	\item кроссплатформенность.
\end{itemize}

Особое значение для разрабатываемой системы имеет механизм сигналов и слотов, обеспечивающий удобное взаимодействие между независимыми компонентами приложения без жёсткой связанности между ними.

Дополнительным преимуществом PySide6 является наличие большого числа готовых интерфейсных компонентов, что сокращает объём низкоуровневой реализации и упрощает сопровождение программного кода.

\subsubsection{Библиотека OpenCV}

Для выполнения операций компьютерного зрения была использована библиотека OpenCV.

OpenCV является одной из наиболее распространённых библиотек для обработки изображений и предоставляет широкий набор инструментов для анализа графических данных.

В рамках разработанной системы библиотека используется для решения следующих задач:

\begin{itemize}
	\item преобразование форматов изображений;
	\item построение ограничивающих рамок;
	\item нанесение текстовых меток;
	\item сохранение обработанных изображений;
	\item цветовые преобразования;
	\item подготовка изображений к последующей визуализации.
\end{itemize}

Выбор OpenCV обусловлен следующими преимуществами:

\begin{itemize}
	\item высокой производительностью;
	\item оптимизацией под работу с большими изображениями;
	\item поддержкой множества форматов данных;
	\item эффективной интеграцией с NumPy;
	\item совместимостью с нейросетевыми моделями.
\end{itemize}

\subsubsection{Библиотека NumPy}

Для эффективной работы с массивами данных и преобразования изображений использовалась библиотека NumPy.

Данная библиотека обеспечивает высокопроизводительные операции над многомерными массивами и широко применяется в задачах обработки изображений и машинного обучения.

В разработанной системе NumPy используется для:

\begin{itemize}
	\item преобразования буферов изображений;
	\item подготовки данных для нейросетевой модели;
	\item изменения структуры массивов;
	\item оптимизации операций обработки изображений.
\end{itemize}

Использование NumPy обеспечивает эффективное взаимодействие между компонентами PySide6, OpenCV и нейросетевой моделью.

\subsubsection{Семейство нейросетевых моделей YOLOv8}

В качестве алгоритма обнаружения дефектов дорожного покрытия была выбрано семейство моделей YOLOv8.

Подход YOLO основан на одноэтапном обнаружении объектов, она выполняет локализацию и классификацию одновременно за один проход.

В рамках проекта использовалась облегчённая версия модели, обеспечивающая снижение вычислительных затрат при сохранении приемлемого качества обнаружения.

Для адаптации модели к задаче обнаружения дефектов дорожного покрытия было выполнено её дополнительное обучение на специализированном датасете изображений дорожного полотна.

\paragraph{Выбор модели семейства}

При выборе архитектуры нейросетевой модели учитывались следующие требования:

\begin{itemize}
	\item высокая скорость обработки изображений;
	\item достаточная точность обнаружения дефектов;
	\item умеренные требования к вычислительным ресурсам;
	\item простота интеграции в Python-приложение.
\end{itemize}

В рамках семейства YOLOv8 доступно несколько моделей:

\begin{itemize}
	\item YOLOv8n (nano);
	\item YOLOv8s (small);
	\item YOLOv8m (medium);
	\item YOLOv8l (large);
	\item YOLOv8x (extra large).
\end{itemize}

Для настоящего проекта была выбрана облегчённая версия YOLOv8n. Это обеспечивает существенное сокращение времени обработки.

Архитектурно модель включает следующие основные компоненты:

\begin{itemize}
	\item backbone-сеть для извлечения признаков;
	\item neck-блок для объединения признаков различных масштабов;
	\item detection head для локализации объектов и определения классов.
\end{itemize}

Выбор данной модификации обусловлен следующими причинами.

\begin{enumerate}
	\item Модель обладает минимальным количеством параметров, что снижает требования к вычислительным ресурсам.

	\item YOLOv8n обеспечивает высокую скорость инференса, что особенно важно при пакетной обработке большого количества изображений дорожного покрытия.

	\item Облегчённая архитектура позволяет использовать систему на персональных компьютерах средней производительности без необходимости применения специализированных GPU-решений.
\end{enumerate}

Несмотря на компактность, YOLOv8n демонстрирует достаточную точность для задач обнаружения локальных дефектов дорожного покрытия при условии корректной подготовки обучающей выборки.

Дополнительным преимуществом выбора YOLOv8 является наличие библиотеки \texttt{Ultralytics}, предоставляющей удобный интерфейс для:

\begin{itemize}
	\item обучения модели;
	\item выполнения инференса;
	\item сохранения весов;
	\item настройки параметров обнаружения;
	\item интеграции в пользовательские приложения.
\end{itemize}

\subsection{Внедрение нейросетевой модели}

\subsubsection{Структура YOLOv8n}

Структурно модель состоит из трёх основных блоков: Backbone, Neck и Detection Head. Её упрощенная схема представлена на рисунке \ref{yolov8:image}.

\paragraph{Backbone}
Backbone (на схеме — область со слоями) построен на модифицированной CSPDarknet. В нём используются блоки \texttt{C2f} (CSPNet с двумя свёртками, объединяющий выходы нескольких внутренних блоков для лучшего распространения информации) и модуль \texttt{SPPF} (быстрое пирамидальное объединение через max pooling, расширяющее рецептивное поле). Задача backbone — извлечь из входного изображения ($640\times640\times3$) карты признаков, описывающие текстуру, границы и форму дефектов.

\paragraph{Neck}
Neck реализован как \texttt{PAN-FPN} (Path Aggregation Network + Feature Pyramid Network). На схеме он представлен блоком «Feature fusion». Сюда поступают признаки от разных слоёв backbone — сплошными стрелками передаются глубокие семантические признаки, пунктирными — детальные пространственные признаки. Neck объединяет их с помощью повышения разрешения (upsampling) и конкатенации. Это формирует многомасштабное представление, необходимое для обнаружения как мелких трещин, так и крупных выбоин.

\paragraph{Detection Head}
Голова (Head) на схеме разделена на две ветви: «Classification» (классификация — определение типа дефекта) и «Regression» (регрессия — уточнение координат ограничивающей рамки). Такая архитектура называется decoupled head (раздельная голова). Модель является anchor-free (не использует заранее заданные якоря), что упрощает обучение. Для повышения точности локализации применяется Distribution Focal Loss (DFL) — функция потерь, которая предсказывает распределение вероятностей для границ бокса, а не единственное значение. На выходе модель формирует три тензора для масштабов $80\times80$ (мелкие объекты), $40\times40$ (средние) и $20\times20$ (крупные).


\begin{figure}[H]
	\center{\includegraphics[width=1\linewidth]{yolov8}}
	\caption{Упрощеная схема структуры модели нейросети}
	\label{yolov8:image}
\end{figure}



\subsubsection{Процесс обучения модели}

Эффективность работы системы автоматического обнаружения дефектов дорожного покрытия напрямую зависит от качества обучения нейросетевой модели. Для адаптации выбранной архитектуры к предметной области был выполнен этап дообучения предварительно обученной модели на специализированном наборе изображений дорожного покрытия.

В качестве исходной модели использовалась предварительно обученная версия \texttt{YOLOv8n}, содержащая веса, полученные на крупномасштабном датасете общего назначения.

Использование предварительно обученной модели позволяет существенно сократить время обучения и повысить итоговую точность обнаружения за счёт использования уже сформированных низкоуровневых признаковых представлений.

Данный подход особенно эффективен для задач, где объём специализированной обучающей выборки ограничен.

Процесс обучения выполнялся с использованием отдельного программного модуля \texttt{train\_model.py}.

Инициализация модели осуществлялась следующим образом:

\begin{lstlisting}[language=Python]
	model = YOLO("models/yolov8n.pt")
\end{lstlisting}

На данном этапе происходила загрузка базовой архитектуры YOLOv8n с предобученными весами.

После инициализации запускался процесс обучения:

\begin{lstlisting}[language=Python]
	model.train(
	data="dataset/pothole_dataset/data.yaml",
	epochs=50,
	batch=16,
	device=0
	)
\end{lstlisting}

Обучение выполнялось на специализированном датасете дорожных дефектов, содержащем размеченные изображения участков дорожного покрытия.

Каждое изображение в обучающей выборке сопровождалось аннотацией, содержащей координаты ограничивающих прямоугольников, локализующих области дефектов.

Структура набора данных соответствовала формату YOLO для обеспечения совместимости датасета с инструментами библиотеки Ultralytics и включала следующие компоненты:

\begin{itemize}
	\item каталог обучающих изображений;
	\item каталог валидационных изображений;
	\item файлы аннотаций;
	\item конфигурационный файл \texttt{data.yaml}.
\end{itemize}

Файл конфигурации содержал информацию о расположении выборок, количестве классов и их наименованиях.

Процесс обучения включал несколько последовательных этапов.

На первом этапе выполнялась загрузка обучающих данных и их предварительная обработка.

Предварительная обработка включала:

\begin{itemize}
	\item нормализацию изображений;
	\item масштабирование;
	\item преобразование координат аннотаций;
	\item формирование обучающих батчей.
\end{itemize}

На втором этапе происходило прямое распространение сигнала через нейронную сеть.

В процессе прямого прохода модель формировала прогнозы относительно наличия дефектов, их координат и уровня достоверности обнаружения.

На третьем этапе вычислялась функция потерь, включающая несколько составляющих:

\begin{itemize}
	\item ошибку локализации;
	\item ошибку классификации;
	\item ошибку оценки достоверности.
\end{itemize}

Полученное значение ошибки использовалось для обратного распространения градиента и корректировки весовых коэффициентов сети.

Данный процесс повторялся для всех батчей текущей эпохи.

После завершения каждой эпохи выполнялась валидация модели на отдельной выборке.

Валидация позволяла контролировать степень обобщающей способности модели и предотвращать переобучение.

В ходе обучения автоматически сохранялись контрольные точки модели.

Наиболее эффективная версия весов сохранялась в файл:

\begin{lstlisting}[language=Python]
	runs/detect/train/weights/best.pt
\end{lstlisting}

Именно данный файл впоследствии использовался в основной программной системе для выполнения обнаружения дефектов.

Использование отдельного файла весов обеспечивает разделение этапов обучения и эксплуатации модели.

Это позволяет:

\begin{itemize}
	\item переобучать модель независимо от основной системы;
	\item использовать различные версии модели;
	\item при необходимости быстро выполнять замену весов.
\end{itemize}

Применение такого подхода повышает гибкость программной системы и упрощает её дальнейшую модернизацию.

В результате обучения была получена специализированная модель, адаптированная к задаче обнаружения дефектов дорожного покрытия и пригодная для интеграции в пользовательское приложение.

\subsubsection{Использование обученной модели в программной системе}

После завершения процесса обучения полученная нейросетевая модель была интегрирована в разработанную программно-информационную систему для выполнения автоматического обнаружения дефектов дорожного покрытия на пользовательских изображениях.

Интеграция модели реализована посредством подключения файла весовых коэффициентов, полученных в процессе обучения.

Загрузка обученной модели выполняется при инициализации основного окна приложения:

\begin{lstlisting}[language=Python]
	self.detector = Detector(
	YOLO("runs/detect/train/weights/best.pt"),
	conf_threshold=0.25
	)
\end{lstlisting}

\subsection{Программная реализация системы}

Система обнаружения дефектов дорожного покрытия построена на основе модульной архитектуры, в которой каждый компонент выполняет строго определённую функциональную роль. Все модули взаимодействуют через центральный управляющий класс, обеспечивающий координацию процессов обработки изображений и работы пользовательского интерфейса.

Система включает в себя следующие модули:

\begin{enumerate}
	\item Главный модуль приложения: обеспечивает координацию всех подсистем и управление состоянием интерфейса.  
	\item Модуль обнаружения объектов: реализует обработку изображений с использованием нейросетевой модели YOLOv8.
	\item Модуль просмотра изображений: обеспечивает просмотр входных и обработанных изображений в отдельном окне.
	\item Модуль работы с файлами: отвечает за выбор, сохранение и загрузку данных, а также формирование отчётов.
	\item Модуль галереи результатов: реализует отображение результатов обработки в виде сеточной структуры.
	\item Модуль управления визуальным оформлением: обеспечивает переключение между светлой и тёмной темой оформления.
\end{enumerate}

На рисунке \ref{DiagramComp:image} представлена диаграмма компонентов проектируемой системы.

\begin{figure}[H]
	\center{\includegraphics[width=0.8\linewidth]{DiagramComp}}
	\caption{Диаграмма компонентов проектируемой системы}
	\label{DiagramComp:image}
\end{figure}

\subsubsection{Главный модуль приложения}

Главный модуль является центральным управляющим компонентом системы. Он обеспечивает инициализацию всех подсистем и координацию их взаимодействия.

В задачи главного модуля входят:

\begin{itemize}
	\item инициализация графического интерфейса пользователя;
	\item создание вспомогательных модулей системы;
	\item управление состоянием приложения;
	\item запуск процесса обработки изображений;
	\item отображение статистики работы системы;
	\item обработка пользовательских событий.
\end{itemize}

Главный модуль агрегирует все функциональные компоненты системы и обеспечивает их согласованную работу в рамках единого приложения.

\subsubsection{Модуль обнаружения объектов}

Модуль обнаружения объектов предназначен для обработки входных изображений с использованием обученной нейросетевой модели YOLOv8n.

Модуль выполняет последовательную обработку изображений, включающую подготовку данных, выполнение нейросетевого анализа и визуализацию результатов. В процессе работы осуществляется преобразование изображений в формат, совместимый с моделью, после чего выполняется инференс и получение координат обнаруженных объектов.

Результаты детекции визуализируются в виде ограничивающих прямоугольников с указанием класса объекта и уровня уверенности модели. Для оптимизации производительности предусмотрено ограничение максимального размера обрабатываемых изображений.

\subsubsection{Модуль просмотра изображений}

Модуль просмотра изображений предназначен для последовательного просмотра отдельных изображений в увеличенном масштабе.

Модуль обеспечивает:

\begin{itemize}
	\item отображение изображения в адаптивном масштабе;
	\item навигацию между изображениями;
	\item автоматическое масштабирование под размер окна.
\end{itemize}

Данный компонент позволяет пользователю детально анализировать результаты работы системы.


\subsubsection{Модуль работы с файлами}

Модуль работы с файлами обеспечивает взаимодействие программы с файловой структурой операционной системы.

Основные функции модуля включают:

\begin{itemize}
	\item выбор входных изображений;
	\item сохранение результатов обработки;
	\item формирование текстового отчёта;
	\item загрузку ранее сохранённых результатов;
	\item управление временными файлами обработки.
\end{itemize}

В процессе сохранения данных формируется отчёт, содержащий информацию о количестве обработанных изображений, числе обнаруженных объектов, среднем количестве детекций на изображение и времени обработки.

\subsubsection{Модуль галереи изображений}

Модуль галереи обеспечивает визуальное представление как исходных, так и обработанных изображений, а также отвечает за перестроение интерфейса при изменении состояния приложения.

Основные функции модуля:

\begin{itemize}
	\item отображение миниатюр входных изображений;
	\item формирование миниатюр обработанных изображений;
	\item размещение элементов в сетке;
	\item отображение количества обнаруженных объектов для каждого изображения;
	\item адаптацию интерфейса к изменению размеров окна.
\end{itemize}

Галерея изображений обеспечивает удобный визуальный анализ входных и обработанных данных.

\subsubsection{Модуль управления визуальным оформлением}

Модуль управления темой отвечает за изменение визуального оформления пользовательского интерфейса.

Система поддерживает два режима отображения: светлый и тёмный. Выбранная тема сохраняется в настройках приложения и автоматически применяется при следующем запуске системы.

Модуль изменяет стилизацию основных элементов интерфейса, включая окна, кнопки, панели и элементы управления.


\subsection{Форматы хранения и обработки данных}

В процессе функционирования программно-информационной системы используются несколько категорий данных, обеспечивающих полный цикл обработки изображений дорожного покрытия: входные изображения, временные результаты обработки, итоговые обработанные изображения и текстовый отчёт.

Организация хранения данных построена таким образом, чтобы обеспечить удобство доступа, минимизацию избыточности и возможность последующего анализа результатов обнаружения дефектов.

\subsubsection{Входные изображения}

Входными данными системы являются цифровые изображения дорожного покрытия, выбираемые пользователем через графический интерфейс программы.

Поддерживаются следующие графические форматы:

\begin{itemize}
	\item JPEG (*.jpg, *.jpeg);
	\item PNG (*.png);
	\item BMP (*.bmp).
\end{itemize}

Выбор данных форматов обусловлен их широкой распространённостью и поддержкой библиотек Qt и OpenCV.

При загрузке изображений осуществляется автоматическая проверка корректности формата и возможность чтения файла. Для корректного отображения ориентации применяется автоматическое преобразование метаданных изображения посредством метода \texttt{setAutoTransform()}.

При необходимости изображения масштабируются до допустимого размера без нарушения пропорций, что позволяет снизить вычислительную нагрузку при обработке.

\subsubsection{Временные результаты}

Промежуточные результаты работы сохраняются во временную директорию, определяемую пользователем либо устанавливаемую по умолчанию.

Временная папка содержит изображения после обработки нейронной сетью.

При запуске новой обработки временная директория автоматически пересоздаётся, что исключает накопление устаревших результатов.

Путь к директории сохраняется в пользовательских настройках с использованием механизма \texttt{QSettings}.

\subsubsection{Итоговые изображения}

Итоговые результаты сохраняются в формате JPEG.

Каждое обработанное изображение содержит:

\begin{itemize}
	\item исходное изображение дорожного покрытия;
	\item прямоугольные ограничивающие рамки обнаруженных дефектов;
	\item подпись класса;
	\item числовое значение достоверности обнаружения.
\end{itemize}

Имена файлов формируются автоматически по шаблону \texttt{result\_XXXX.jpg}, где \texttt{XXXX} --- порядковый номер изображения.

\subsubsection{Текстовый отчёт}

По завершении обработки формируется текстовый отчёт в формате TXT.

Отчёт содержит:

\begin{itemize}
	\item количество загруженных изображений;
	\item количество успешно обработанных изображений;
	\item суммарное число обнаруженных дефектов;
	\item среднее количество дефектов на изображение;
	\item общее время обработки;
	\item детализацию по каждому изображению.
\end{itemize}

Структура отчёта обеспечивает возможность последующего документирования результатов и использования данных в аналитических целях. Также он используется системой для восстановления статистической информации о ранее обработанных изображениях при их повторной загрузке.

\subsection{Надёжность и производительность}

Одним из ключевых требований к разрабатываемой системе является обеспечение стабильной работы при обработке больших наборов изображений, а также эффективное использование вычислительных ресурсов при выполнении детекции.

Для достижения требуемого уровня надёжности и производительности в системе реализован комплекс решений, связанных с обработкой исключительных ситуаций, оптимизацией загрузки данных и рациональным управлением памятью.

\subsubsection{Обработка ошибок}

Система предусматривает обработку основных исключительных ситуаций, возникающих при работе с файлами, пользовательским вводом и нейросетевой обработкой.

К основным контролируемым ошибкам относятся:

\begin{itemize}
	\item попытка открытия неподдерживаемого файла;
	\item невозможность чтения изображения;
	\item отсутствие доступа к директории сохранения;
	\item повреждение файлов результатов;
	\item прерывание обработки пользователем.
\end{itemize}

Пользователь уведомляется о возникших проблемах посредством диалоговых окон и сообщений в строке состояния.

\subsubsection{Оптимизация загрузки изображений}

Для повышения скорости работы применяется предварительное масштабирование изображений на этапе чтения.

Использование класса \texttt{QImageReader} позволяет загружать изображения сразу в уменьшенном разрешении, что существенно снижает:

\begin{itemize}
	\item время декодирования;
	\item объём используемой оперативной памяти;
	\item нагрузку на процессор.
\end{itemize}

Дополнительно реализована отложенная перестройка сетки результатов посредством таймера, что предотвращает избыточные операции перерасчёта интерфейса.

\subsubsection{Управление памятью}

Рациональное использование памяти достигается за счёт:

\begin{itemize}
	\item последовательной обработки изображений;
	\item удаления устаревших временных объектов;
	\item использования масштабированных копий вместо полноразмерных изображений;
	\item применения механизма \texttt{deleteLater()} для корректного освобождения ресурсов Qt.
\end{itemize}

При отображении миниатюр используются только уменьшенные копии изображений, что позволяет сохранять отзывчивость интерфейса.

\subsubsection{Производительность обработки}

Основное время вычислений приходится на выполнение инференса модели YOLOv8.

Общая скорость обработки зависит от:

\begin{itemize}
	\item разрешения входных изображений;
	\item количества изображений;
	\item вычислительной мощности GPU;
	\item выбранного порога обнаружения.
\end{itemize}

Для контроля производительности реализован замер общего времени обработки.

Полученные показатели отображаются пользователю в интерфейсе программы и фиксируются в отчёте.

\subsection{Проектирование пользовательского интерфейса}

Проектирование пользовательского интерфейса разработанной системы осуществлялось с учётом требований, представленных в пункте 2.4 технического задания.

Интерфейс системы реализован на основе графического фреймворка PySide6 и представляет собой настольное приложение с событийно-ориентированной моделью взаимодействия.

Интерфейс программно-информационной системы включает следующие основные элементы:

\begin{itemize}
	\item главное окно приложения;
	\item меню управления обработкой изображений;
	\item область отображения входных изображений;
	\item область отображения результатов обработки;
	\item окно расширенного просмотра изображений;
	\item панель статистики обработки;
	\item меню управления файловыми операциями;
	\item меню изменения вида главного окна и сортировки результатов.
\end{itemize}

На рисунке \ref{Maket:image} представлен макет пользовательского интерфейса главного окна программной системы.

\begin{figure}[H]
	\center{\includegraphics[width=1\linewidth]{Maket}}
	\caption{Макет пользовательского интерфейса главного окна}
	\label{Maket:image}
\end{figure}

\subsection{Верификация технического проекта}

Проверка работы программно-информационной системы должна включать:

\begin{enumerate}
	\item Запуск приложения и корректную инициализацию главного окна.
	\item Загрузку изображений поддерживаемых форматов.
	\item Проверку обработки изображений нейросетевой моделью YOLOv8n.
	\item Корректное отображение исходных и обработанных изображений в галереях.
	\item Открытие изображений в режиме полноразмерного просмотра.
	\item Навигацию между изображениями в модуле просмотра.
	\item Сохранение обработанных изображений в выбранную директорию.
	\item Формирование текстового отчёта по результатам обработки.
	\item Загрузку ранее сохранённых результатов.
	\item Переключение между светлой и тёмной темами оформления.
	\item Проверку сохранения пользовательских настроек после перезапуска приложения.
	\item Обработку изображений различного разрешения.
	\item Корректную очистку временных файлов при повторной обработке.
	\item Устойчивую работу системы при последовательной обработке большого количества изображений.
	\item Корректное завершение работы приложения с освобождением используемых ресурсов.
\end{enumerate}